\renewcommand{\Anno}{1998}

\renewcommand{\NumeroProblema}{1}

\renewcommand{\Testo}{Determinare per quali valori del parametro reale
$\alpha$ la seguente equazione
$$1998^{\,|\sin x|}=|\sin(\alpha x)|^{1998}$$
ha almeno una soluzione reale.}

\renewcommand{\Risposta}{L'equazione ha almeno una soluzione reale se
e solo se $\alpha$ è un numero razionale che, scritto sotto forma
di frazione irriducibile, ha denominatore pari.}

\renewcommand{\Aiutino}{Quanto vale, come minimo, il termine di
sinistra?  Quanto vale, come massimo, il termine di destra?  Quando
possono essere uguali?}

\renewcommand{\SoluzioneUfficiale}{
}

\renewcommand{\SoluzioniAlternative}{
}

\renewcommand{\Osservazioni}{
}

\renewcommand{\Commenti}{
}
% Aggiunta redazionale: la soluzione preesistente non viene modificata.
\providecommand{\SoluzioneRedazionale}{}
\renewcommand{\SoluzioneRedazionale}{
% SSSUP-BOZZA-AI: SoluzioneRedazionale; autore=ChatGPT; data=2026-11-08; stato=DA-VERIFICARE
% SSSUP-PROVENIENZA: AI: elaborazione nuova dal testo
Il primo membro è almeno $1$, mentre il secondo è al più $1$. L'uguaglianza richiede quindi simultaneamente
\[
\sin x=0,\qquad |\sin(\alpha x)|=1.
\]
Esistono pertanto interi $k,m$, con $k\neq0$, tali che
\[
x=k\pi,\qquad \alpha k\pi=\frac{(2m+1)\pi}{2}.
\]
Ne segue $\alpha=(2m+1)/(2k)$: riducendo la frazione, il denominatore resta pari, perché il numeratore è dispari.

Viceversa, sia $\alpha=p/q$ una frazione irriducibile con $q>0$ pari. Allora $p$ è dispari e basta prendere $x=q\pi/2$: si ha $\sin x=0$ e $|\sin(\alpha x)|=|\sin(p\pi/2)|=1$. La condizione richiesta è dunque esattamente quella indicata nella risposta preesistente.
}
